$\hat g_\lambda$ is an isometry $(T, h_{u+\lambda}) \to (T, h_u)$: it preserves $G$ on covector pairs (Lemma \claimref{sheaf.base-unipotents-wall}), is the identity on the adjoined plane, and on the coupling $h_u(\hat g\alpha, \nO) = \nu(g\alpha)(u) = \nu\alpha(u) + \nu\alpha(\lambda)k(u) = \nu\alpha(u+\lambda) = h_{u+\lambda}(\alpha, \nO)$, identically with $\mu$. Composing with $\Phi_\lambda^{-1} : (T, h_u) \to (T, h_{u+\lambda})$ gives the automorphism; its restriction to $T^*M$ is $g_\lambda$ since $\Phi_\lambda$ is the identity there, both factors fix $k$, and an isometry of $(T, h_u)$ mapping $T^*M$ to itself preserves the grading (it is induced by a map of the generating module) and the subalgebra $T^*M$ generates. The contragredient of $g_\lambda = \id + k \otimes \lambda$ (the map $\alpha \mapsto \alpha + \alpha(\lambda)k$) is $(g_\lambda^{-1})^{\mathsf{T}} = \id - \lambda \otimes k$, that is $u' \mapsto u' - k(u')\lambda$, which is Lemma \claimref{sheaf.base-unipotents-wall}'s translation on candidates. A fixed point would satisfy $u' = u' - \lambda$ for every $\lambda$. The fixed set: write $x = a + b\nO + c\nInf$ with $a \in T^*M$, and let $\tilde\lambda$ be the normalised representative of $\lambda$. Then $\Phi_\lambda^{-1}x = x - (b\nu + c\mu)\tilde\lambda$, and applying $\hat g_\lambda$ adds $[a(\lambda) - (b\nu + c\mu)\tilde\lambda(\lambda)]\,k$; so $\theta_\lambda x = x$ iff $(b\nu + c\mu)\tilde\lambda \in \R k$ and the bracket vanishes. For $\lambda \neq 0$, $\tilde\lambda \notin \R k$ (it pairs to $\lambda \neq 0$ under $G$, while $k$ pairs to zero), so $b\nu + c\mu = 0$, and $h_u(k, x) = G(k, a) + (b\nu + c\mu)k(u) = 0$. Injectivity of $v \mapsto \tilde v$ is the solve of Lemma \claimref{sheaf.solve-closed-radical-gauge}: $\tilde v$ determines $v$ through $\alpha(v) = G(\alpha, \tilde v)$.
